
\section{A sufficient condition from the squared momentum equation}
\label{r:sec:first-rectangle}

The temporal and viscous terms can be separated directly in an integrated
identity. Write
\[
 \mathcal G=\Pp f-\Pp[(u\cdot\nabla)u],
 \qquad u_t-\nu\Delta u=\mathcal G.
\]
For a strict classical horizon \(S<T\), put
\[
 X(S)=\int_0^S(\norm{u_t}_2^2+\nu^2\norm{\Delta u}_2^2)\dd t,
 \qquad G_0=\norm{\nabla u_0}_2^2.
\]
The cross term in the squared equation is an endpoint energy:
\begin{equation}
 X(S)+\nu\norm{\nabla u(S)}_2^2
 =\nu G_0+\int_0^S\norm{\mathcal G}_2^2\dd t.
 \label{r:eq:first-rectangle-exact}
\end{equation}
Indeed, \(-2\ip{u_t}{\Delta u}=\partial_t\norm{\nabla u}_2^2\).
This identity retains the full projected nonlinear term.

Set
\[
 U=\norm{u_0}_2+\int_0^T\norm f_2\dd t,\qquad
 L=TU^2+U^2/\nu,\qquad
 \Phi=\int_0^T\norm{\Pp f}_2^2\dd t.
\]
The kinetic balance gives \(\sup_t\norm u_2\le U\) and
\(\nu\int_0^S\norm{\nabla u}_2^2\dd t\le U^2/2\).
Using the Bessel Sobolev norm,
\[
 \norm u_{H^2}^2=\norm u_2^2+2\norm{\nabla u}_2^2+\norm{\Delta u}_2^2.
\]
Consequently the first rectangle
\(R_{0,1}(S)=\int_0^S(\norm u_{H^2}^2+\norm{u_t}_2^2)\dd t\)
satisfies
\begin{equation}
 \begin{aligned}
 R_{0,1}(S)&\le L+\max(1,\nu^{-2})X(S),\\
 X(S)&\le\max(1,\nu^2)R_{0,1}(S).
 \end{aligned}
 \label{r:eq:first-rectangle-equivalence}
\end{equation}
A bound for the full \(\mathcal G\) integral thus supplies both
velocity coordinates without first assuming that rectangle.

The gradient trace identity and Young's inequality also give
\[
 \sup_{t\le S}\norm{\nabla u(t)}_2^2\le G_0+X(S)/\nu,
 \qquad
 \int_0^S\norm u_{H^2}^2\dd t\le L+X(S)/\nu^2.
\]
Fourier inversion gives
\(\norm u_\infty^2\le\norm u_{H^2}^2/(8\pi)\), proved explicitly in
\eqref{r:eq:embedding}. Since
\(\norm{(u\cdot\nabla)u}_2\le\norm u_\infty\norm{\nabla u}_2\),
substitution into \eqref{r:eq:first-rectangle-exact} yields
\begin{equation}
 X(S)\le\nu G_0+2\Phi+
 \frac{(G_0+X(S)/\nu)(L+X(S)/\nu^2)}{4\pi}.
 \label{r:eq:first-rectangle-polynomial}
\end{equation}
This proves a quantified criterion directly from the data. Define
\[
 a=\frac1{4\pi\nu^3},\qquad
 b=\frac{G_0/\nu^2+L/\nu}{4\pi},\qquad
 c=\nu G_0+2\Phi+\frac{G_0L}{4\pi}.
\]
If \(b<1\) and \((1-b)^2>4ac\), then
\begin{equation}
 \sup_{S<T}X(S)\le
 \frac{1-b-\sqrt{(1-b)^2-4ac}}{2a}.
 \label{r:eq:first-rectangle-root}
\end{equation}
To see this, rewrite \eqref{r:eq:first-rectangle-polynomial} as
\(aX^2+(b-1)X+c\ge0\). Its two roots enclose a forbidden interval.
The continuous function \(X\), starting from zero, cannot cross that interval.
For rest datum the criterion becomes
\[
 \frac{U^2(T+\nu^{-1})}{4\pi\nu}
 +\sqrt{\frac{2\Phi}{\pi\nu^3}}<1.
\]
The bound is uniform through \(T\). Without the stated root separation,
\eqref{r:eq:first-rectangle-polynomial} does not supply an upper bound.
